\documentclass[14pt,usepdftitle=false,aspectratio=169]{beamer}
\usepackage{preambule}
\setbeamercolor{structure}{fg=black}
\usepackage{polynomes,al,usuelles,topologie,structures}
\hypersetup{pdftitle={Théorie spectrale -- Transformée de Gelfand}}
\title{Théorie spectrale\\\emph{Transformée de Gelfand}}
\author{}
\date{}
\begin{document}
\begin{frame}
    \titlepage
\end{frame}
\slideq{Propriétés de $A/I$ pour $A$ une algèbre de Banach et $I$ un idéal fermé de $A$}{1/11}
\slider{$A/I$ muni de la norme quotient $\nrm{\pi\l a\r}=\inf[x\in I]{\nrm{a-x}}$ est une algèbre de Banach}{1/11}
\slideq{Propriétés des éléments de $\sigma\l A\r$}{2/11}
\slider{Tout $\omega\in\sigma\l A\r$ vérifie:\linebreak$\omega\l a\r\in\sigma\l a\r$\linebreak$\nnrm\omega=1$\linebreak$\omega$ est une forme linéaire continue}{2/11}
\slideq{Théorème de Gelfand}{3/11}
\slider{Toute algèbre de Banach qui est aussi un corps est isomorphe à $\bbC$\linebreak Toute $\bbR$-algèbre de Banach qui est aussi un corps est isomorphe à $\bbR$, $\bbC$ ou $\bbH$}{3/11}
\slideq{Correspondance des idéaux maximaux d'une algèbre de Banach}{4/11}
\slider{Si $A$ est commutative, pour tout idéal maximal $I$ de $A$, il existe un unique $\omega\in\sigma\l A\r$ tel que $\ker\omega=I$}{4/11}
\slideq{Topologie sur $\sigma\l A\r$}{5/11}
\slider{Topologie de la convergence simple\linebreak Restriction de la topologie faible-$*$ de $A'$}{5/11}
\slideq{Propriété de $\appl{\delta}{X}{\sigma\l\calC\l X\r\r}{x}{\delta_x=\l f\mapsto f\l x\r\r}$ pour $X$ compact}{6/11}
\slider{C'est un homéomorphisme}{6/11}
\slideq{Propriété des idéaux maximaux des algèbres de Banach}{7/11}
\slider{Les idéaux maximaux des algèbres de Banach sont fermés}{7/11}
\slideq{Propriétés topologiques de $\sigma\l A\r$}{8/11}
\slider{C'est un espace compact}{8/11}
\slideq{Propriétés de la transformée de Gelfand}{9/11}
\slider{$\nnrm{\varPhi_A}\leqslant1$\linebreak$\sigma_A\l a\r=\sigma\l\varPhi_A\l a\r\r$\linebreak$\nrm{\varPhi_A\l a\r}=r\l a\r$}{9/11}
\slideq{Spectre de Gelfand d'une $\bbC$-algèbre $A$}{10/11}
\slider{$\sigma\l A\r$ est l'ensemble des morphsimes d'algèbres $A\to\bbC$}{10/11}
\slideq{Transformée de Gelfand de $A$}{11/11}
\slider{$\appl{\varPhi_A}{A}{\calC\l\sigma\l A\r\r}{x}{\l\omega\mapsto\omega\l x\r\r}$\linebreak C'est un morphisme d'algèbres}{11/11}
\end{document}